\section{Inférence sur les paramètres}
\label{sec:InférenceParamètres}

On a montré comment obtenir des estimations pour les différents paramètres du modèle de régression linéaire simple. Généralement, on souhaite également faire de l'inférence sur les vraies valeurs des paramètres $\beta_0$, $\beta_1$ et $\sigma^2$. Pour ce faire, on utilise la distribution des estimateurs.

\subsection{Distribution des estimateurs}
\label{sec:LoisBetas}

\subsubsection{Distribution de $\widehat{\beta}_1$}
\label{sec:LoiDeHatBeta1}

On note d'abord que $\widehat{\beta}_1$ est une combinaison linéaire des $Y_i$. En effet,
\begin{align*} 
\widehat{\beta}_1 &= \frac{\sum\limits_{i=1}^n (x_i-\overline{x})Y_i}{\sum\limits_{i=1}^n (x_i-\overline{x})^2} = \sum_{i=1}^n a_i Y_i, \cr
\intertext{où}
a_i &= \frac{x_i-\overline{x}}{\sum\limits_{i=1}^n(x_i-\overline{x})^2} 
    =  \frac{x_i-\overline{x}}{S_{xx}}.
\end{align*}
Étant donné qu'une combinaison linéaire de variables aléatoires normales suit également une distribution normale, il s'ensuit que $\widehat\beta_1$ est distribué selon une loi normale. Il faut déterminer l'espérance et la variance de $\widehat\beta_1$. Notons d'abord que
 \begin{eqnarray*}
 \sum_{i=1}^n a_i & = & \frac{\sum\limits_{i=1}^n(x_i-\overline{x})} {S_{xx}} = 0, \\
\sum_{i=1}^n a_i x_i & = & \frac{\sum\limits_{i=1}^n(x_i-\overline{x})x_i } {S_{xx}} =\frac{\sum\limits_{i=1}^n(x_i-\overline{x})(x_i-\overline{x})}{S_{xx}}=\frac{S_{xx}}{S_{xx}}=1.
\end{eqnarray*}
Calculons l'espérance  de $\beta_1$,
\begin{eqnarray*}
 \esp\left(\widehat{\beta}_1\right) &=& \sum_{i=1}^n a_i \esp\left(Y_i\right)  \\
 &=& \sum_{i=1}^n a_i \left(\beta_0 + \beta_1 x_i\right)\\
 &=& \beta_0 \sum_{i=1}^n a_i + \beta_1 \sum_{i=1}^n a_i x_i\\
 &=& \beta_1.
\end{eqnarray*}
Ainsi, l'estimateur $\widehat{\beta}_1$ est un estimateur sans biais de $\beta_1$ puisque $\esp\left(\widehat{\beta}_1\right)  - \beta_1 = 0$.
Calculons la variance de $\widehat\beta_1$, 
\begin{eqnarray*}
 \var\left(\widehat\beta_1\right)&=&\var\left(\sum_{i=1}^n a_i Y_i\right)\\
 &=& \sum_{i=1}^n a_i ^2 \, \var\left(Y_i\right)\\
 &=& \sum_{i=1}^n \frac{(x_i-\overline{x})^2}{S_{xx}^2}  \, \var\left(Y_i\right)\\
 &=& \frac{1}{S^2_{xx}} \sum_{i=1}^n (x_i-\overline{x})^2  \, \var\left(Y_i\right)\\
 &=& \frac{S_{xx}}{S_{xx}^2} \sigma^2 \\
 &=& \frac{\sigma^2}{S_{xx}}.
\end{eqnarray*}
Ainsi, la distribution de $\widehat\beta_1$ est
\begin{eqnarray}
\label{distbeta1}
 \widehat{\beta}_1 \sim N\left( \beta_1, \frac{\sigma^2}{S_{xx}} \right).
 \label{loibeta1}
\end{eqnarray}

\begin{remarque}
Seulement certaines hypothèses de la RLS ont été nécessaires pour calculer cette distribution. Lesquelles? Et que nous indique la forme de cette distribution sur l'estimation de la pente dans un modèle de régression linéaire simple?
\end{remarque}


\subsubsection{Distribution de $\widehat\beta_0$}
\label{DistributionBeta0}
L'estimateur $\widehat\beta_0$ peut également s'exprimer comme une combinaison linéaire des $Y_i$,
\begin{eqnarray*}
 \widehat{\beta}_0 &=& \overline{Y} - \widehat{\beta}_1 \overline{x} \\
 &=& \sum_{i=1}^n \frac{Y_i}{n} - \overline{x} \sum_{i=1}^n a_i Y_i \\
 &=& \sum_{i=1}^n \left( \frac{1}{n} - a_i \overline{x}\right) Y_i.
\end{eqnarray*}
Par le même raisonnement que pour $\widehat\beta_1$, on a donc que $\widehat\beta_0$ est distribué selon une loi normale. On peut calculer l'espérance et la variance de $\widehat\beta_1$. Pour l'espérance, on a que
\begin{eqnarray*}
 \esp\left(\widehat{\beta}_0\right) &=& \esp\left(\overline{Y} - \widehat{\beta}_1 \overline{x}\right) \\
 &=& \frac{1}{n}\sum_{i=1}^n \esp\left(Y_i\right) - \overline{x}  \,\esp\left(\widehat{\beta}_1\right)  \\
 &=& \frac{1}{n}\sum_{i=1}^n\left(\beta_0+\beta_1 x_i\right)-\overline{x}\beta_1 \\
 & = & \beta_0+\beta_1 \overline x-\beta_1\overline x \\
 & = & \beta_0.
\end{eqnarray*}
Pour la variance, on a
\begin{eqnarray*}
 \var\left(\widehat{\beta}_0\right) &=& \var\left(\overline{Y} - \widehat{\beta}_1 \overline{x}\right) \\
 &=& \var\left(\frac{1}{n}\sum_{i=1}^n Y_i\right)+ \overline{x}^2 \, \var\left(\widehat{\beta}_1\right)-2\overline x \, \cov\left(\overline Y,\widehat\beta_1\right)  \\
 &=& \frac{\sigma^2}{n}+\overline{x}^2\frac{\sigma^2}{S_{xx}} \\
 & = & \sigma^2\left(\frac{1}{n}+\frac{\overline{x}^2}{S_{xx}}\right).
\end{eqnarray*}
Pour obtenir l'expression finale pour cette variance, il faut démontrer que $\cov[\overline Y,\widehat\beta_1]=~0$, ce qui est laissé en exercice. 
On obtient ainsi la loi pour $\widehat\beta_0$,
\begin{eqnarray}
\label{distbeta0}
 \widehat{\beta}_0 \sim \N\left\{\beta_0, \left(\frac{1}{n}+\frac{\overline{x}^2}{S_{xx}}\right)\sigma^2\right\}.
\label{loibeta0}
\end{eqnarray}

\begin{exemple}
Calculons la variance de $\widehat\beta_0$ et $\widehat\beta_1$ de l'exemple de la section \ref{exemple-RLS-interpretation}. On doit pour ce faire connaître la valeur de la variance des termes d'erreur. Supposons ici que $\sigma^2=1$. On sait déjà que $n=10,\overline x=145$ et $S_{xx}=8250$. Par conséquent on a
\begin{align*}
    \var[\widehat\beta_0] &= \left( \frac{1}{10}+\frac{145^2}{8250} \right)(1)=2,648, \cr
    \var[\widehat\beta_1] &= \frac{1^2}{8250}=1,212\times 10^{-4}.
\end{align*}
\end{exemple}


\subsection{Intervalles de confiance pour $\beta_0$ et $\beta_1$}
\label{sec:IntervallesDeConfiance}

On peut utiliser les distributions de $\widehat{\beta}_0$ et $\widehat{\beta}_1$ pour construire des intervalles de confiance pour $\beta_0$ et $\beta_1$. Pour déterminer l'intervalle de confiance pour $\beta_1$, on centre et réduit la variable aléatoire pour obtenir
\begin{equation}
\frac{\widehat{\beta}_1-\beta_1}{\sqrt{\sigma^2/S_{xx}}}\sim \N(0,1).
\label{disbeta0std}
\end{equation}
%équations \eqref{distbeta0} et \eqref{distbeta1})
\vspace{-0.05 in}
Si $\sigma^2$ était connue, un intervalle de confiance de niveau $100(1-\alpha)\%$ serait donné par 
$$\left[\widehat{\beta}_1\pm z_{1-\alpha/2} \frac{\sigma}{\sqrt{S_{xx}}} \right].$$ 
Or, $\sigma^2$ n'est généralement pas connue et on doit l'estimer. L'estimateur du maximum de vraisemblance de cette quantité est 
$$ \widehat{\sigma}_{MV}^2 = \sum_{i=1}^n \frac{\widehat{\varepsilon}_i^2}{n},$$ mais ce dernier est biaisé.
%Montrer que c'est le cas.
On utilisera donc plutôt l'estimateur non biaisé défini par $$\widehat{\sigma}^2=  \frac{1}{n-2} \sum_{i=1}^n \widehat{\varepsilon}_i^2.$$

\vspace{-0.05 in}
Remplacer $\sigma^2$ par $\widehat\sigma^2$ dans le terme de gauche de \eqref{disbeta0std} a toutefois pour conséquence que cette variable aléatoire ne suit plus une loi normale. On peut néanmoins montrer, en utilisant des notions de statistique mathématique, que
\begin{equation}
 \frac{\displaystyle\frac{\widehat{\beta}_1-\beta_1}{\displaystyle\sigma/\sqrt{S_{xx}}}}
 {  \sqrt{\frac{\displaystyle\widehat\sigma^2}{\displaystyle\sigma^2} } } \sim \mathcal{T}_{n-2}.
\label{disbeta0stds}
\end{equation}
On peut alors déduire un intervalle de confiance de niveau $100(1-\alpha)\%$ pour $\beta_1$,
\begin{eqnarray}
\left[\widehat{\beta}_1 \pm t_{n-2; 1-\alpha/2} \left(\frac{\widehat\sigma}{\sqrt{S_{xx}}}\right)\right],
\end{eqnarray}
où $t_{n-2; 1-\alpha/2}$ dénote le $(1-\alpha/2)$-ième quantile d'une loi $\mathcal{T}$ à $(n-2)$ degrés de liberté.\\

 Il est possible de construire de la même manière un intervalle de confiance pour $\beta_0$. On obtient l'intervalle de confiance 
 \begin{eqnarray}
\left[\widehat{\beta}_0\pm t_{n-2; 1-\alpha/2}\; \widehat\sigma\sqrt{{1\over n}+{\overline{x}^2\over S_{xx}}}\right].
\label{ci:beta0}
\end{eqnarray}

\vspace{0.05in}
\begin{remarque}
Pour arriver à \eqref{disbeta0stds}, on note que
\begin{eqnarray*}
 \frac{\frac{\displaystyle\widehat{\beta}_1-\beta_1}{\displaystyle\sigma/\sqrt{S_{xx}}}}{\sqrt{\frac{\displaystyle\sum_{i=1}^n \widehat{\varepsilon}_i^2}{\displaystyle\sigma^2 (n-2)}}}=\frac{U}{\sqrt{V/\nu}},
\end{eqnarray*}
où $U=\left(\widehat{\beta}_1-\beta_1\right) \big/ \left(\sigma/\sqrt{S_{xx}}\right)$ est distribué selon une loi normale centrée réduite et $V= \left(\sum_{i=1}^n \widehat{\varepsilon}_i^2 \right)\big/ \sigma^2$ suit une loi du khi-carré à $\nu=(n-2)$ degrés de liberté, puisque c'est une somme du carré de variables aléatoires normales.
Comme il s'agit d'un rapport entre une variable aléatoire normale centrée réduite et la racine carrée d'une variable khi-carré divisée par son degré de liberté  $(n-2)$, on obtient la distribution d'une loi $\mathcal{T}$ de Student à $(n-2)$ degrés de liberté ($\mathcal{T}_{n-2}$),
\begin{eqnarray*}
 \frac{\frac{\displaystyle\widehat{\beta}_1-\beta_1}{\displaystyle\sigma/\sqrt{S_{xx}}}}{\sqrt{\frac{\displaystyle\sum_{i=1}^n \widehat{\varepsilon}_i^2}{\displaystyle\sigma^2 (n-2)}}} \sim \frac{\N(0,1)}{\sqrt{\chi_{n-2}^2/(n-2)}} \sim \mathcal{T}_{n-2}.
 \end{eqnarray*}
\end{remarque}


\begin{exemple}
\label{exp:int_beta}
Calculons un intervalle de confiance pour $\beta_0$ et $\beta_1$ de l'exemple \ref{exemple-RLS-interpretation}. La plupart des quantités nécessaires ont déjà été calculées. Il ne manque que
$$\widehat\sigma=\sqrt{\frac{1}{10-2} \sum_{i=1}^{10}\widehat\varepsilon_i^2}=\sqrt{\frac{7,224}{8}}=0,950.$$
Si on fixe $\alpha=5\%$, il faut calculer $t_{n-2; 0,975}$. On peut trouver dans une table que ce quantile est  $t_{8;0,975}=2,306$. On obtient alors les intervalles suivants:
$$\left[\widehat{\beta}_0\pm 2,306\times0,950\sqrt{1/10+145^2/8250} \right]=\left[\widehat{\beta}_0\pm 3,566\right]=\left[-6,305;0,827\right],$$
$$\left[\widehat{\beta}_1\pm 2,306\times 0,950/\sqrt{8250}\right]=\left[\widehat{\beta}_1\pm 0,024\right]=[0,459;0,507].$$
\end{exemple}


\subsection{Tests d'hypothèses}

On peut également se baser sur les distributions déterminées à la section \ref{sec:LoisBetas} pour mettre en {\oe}uvre des tests d'hypothèses sur la valeur des coefficients. Considérons le test bilatéral
\begin{align*}
    \mathcal{H}_0&: \; \beta_j = c, \cr
    \mathcal{H}_1&: \; \beta_j \neq c.
\end{align*}
Notons 
\begin{align*}
    \widehat{\mathrm{Var}}\left(\widehat{\beta}_0\right)
    &= \widehat\sigma^2\left({1\over n}+{\overline{x}^2\over S_{xx}}\right), \\
    \widehat{\mathrm{Var}}\left(\widehat{\beta}_1\right)
    &= \frac{\widehat\sigma^2}{S_{xx}}.
\end{align*}
Sous $\mathcal{H}_0$, 
$$\frac{\widehat{\beta}_j - c}{\widehat{\mathrm{Var}}\left(\widehat{\beta}_j\right)} \sim \mathcal{T}_{n-2}.$$
Soit $t_{obs} = \left(\widehat{\beta}_j - c \right) \big/ \widehat{\mathrm{Var}}\left(\widehat{\beta}_j\right)$ la statistique observée du test. La décision du test d'hypothèses au seuil $\alpha$ peut être prise de trois façons.
\begin{enumerate}[label=(\roman*)]
    \item Valeur critique: Soit le quantile théorique $t_{n-2; 1-\alpha/2}$. Si $|t_{obs}| > t_{n-2; 1-\alpha/2}$, l'hypothèse nulle $\mathcal{H}_0: \; \beta_j = c$ est rejetée.
    \item Valeur-p: Soit $\alpha^* = \prob(|t_{n-2}|>|t_{\text{obs}}|) = 2 \prob(t_{n-2}>|t_{\text{obs}}|)$ la valeur-p. Si $\alpha^* < \alpha$,  l'hypothèse nulle $\mathcal{H}_0: \; \beta_j = c$ est rejetée.
    \item Intervalle de confiance: Si la valeur $c$ n'est pas comprise dans l'intervalle de confiance pour $\beta_j$ de niveau $1-\alpha$, l'hypothèse nulle $\mathcal{H}_0: \; \beta_j = c$ est rejetée.
\end{enumerate}
Il est aussi possible de faire des test unilatéraux. Pour le test
\begin{align*}
    \mathcal{H}_0&: \; \beta_j = c, \cr
    \mathcal{H}_1&: \; \beta_j < c,
\end{align*}
les règles de décisions sont
\begin{enumerate}[label=(\roman*)]
    \item Valeur critique: Si $t_{obs} < - t_{n-2; 1-\alpha}$, l'hypothèse nulle $\mathcal{H}_0: \; \beta_j = c$ est rejetée.
    \item Valeur-p: Si $\alpha^* = \prob(t_{n-2}<t_{\text{obs}}) < \alpha$,  l'hypothèse nulle $\mathcal{H}_0: \; \beta_j = c$ est rejetée.
\end{enumerate}
Pour le test
\begin{align*}
    \mathcal{H}_0&: \; \beta_j = c, \cr
    \mathcal{H}_1&: \; \beta_j > c,
\end{align*}
les règles de décisions sont
\begin{enumerate}[label=(\roman*)]
    \item Valeur critique: Si $t_{obs} > t_{n-2; 1-\alpha}$, l'hypothèse nulle $\mathcal{H}_0: \; \beta_j = c$ est rejetée.
    \item Valeur-p: Si $\alpha^* = \prob(t_{n-2}>t_{\text{obs}}) < \alpha$,  l'hypothèse nulle $\mathcal{H}_0: \; \beta_j = c$ est rejetée.
\end{enumerate}

Un cas particulier, souvent d'intérêt, est le test bilatéral avec $c=0$. Pour la pente, ceci correspond aux hypothèses
$\mathcal{H}_0: \beta_1 = 0$, $\mathcal{H}_1: \beta_1 \neq 0$. Si l'hypothèse nulle de ce test est rejetée, le lien linéaire entre $Y$ et $x$ n'est pas significatif et le modèle se réduit à $Y_i = \beta_0 + \varepsilon_i$. Pour l'ordonnée à l'origine, on a $\mathcal{H}_0: \beta_0 = 0$, $\mathcal{H}_1: \beta_0 \neq 0$. Si l'hypothèse nulle de ce test est rejetée, l'ordonnée à l'origine n'est pas significativement différente de 0 et le modèle se réduit à $Y_i = \beta_1 x_i + \varepsilon_i$.


\begin{exemple}
Toujours avec les données de la section \ref{exemple-RLS-interpretation}, on a
$$ t_{obs} = \frac{\widehat{\beta}_1}{\widehat\sigma/\sqrt{S_{xx}}} = \frac{0,483}{0,950/\sqrt{8250}} = 46,17. $$
Un logiciel nous permet d'obtenir $5.35 \times 10^{-11}$ comme valeur de $p$ correspondante. On rejette donc l'hypothèse d'une pente nulle au seuil $\alpha = 0,05$. Notez qu'on aurait pu arriver à la même conclusion simplement à l'aide de l'intervalle de confiance construit à l'exemple \ref{exp:int_beta}. En effet, l'intervalle de confiance à 95\% pour $\beta_1$ ne contient pas la valeur $0$, ce qui permet de rejeter l'hypothèse que $\beta_1 = 0$ au seuil $\alpha = 0,05$.
\end{exemple}


\section{Prévision}
\label{sec:PrévisionRLS}

On a déjà défini $\widehat{Y}_i = \widehat{\beta}_0 + \widehat{\beta}_1 x_i$, la prévision du modèle en $x_i$. Avant de définir un intervalle de confiance, deux intérêts doivent être distingués:
\begin{enumerate}
\item[(i)] l'estimation de $\esp[Y|x=x_0]$, la valeur moyenne de $Y$ pour une valeur $x=x_0$
de la variable explicative.
\item[(ii)] l'estimation de ${Y}_0^*$, soit la valeur d'une observation/r\'ealisation de $Y$ lorsque la variable explicative prend la valeur $x_0$. 
\end{enumerate}
Dans le premier cas, on souhaite estimer la valeur de la droite en $x_0$, alors que dans le deuxième cas, on s'intéresse à une réalisation de la variable réponse, qui n'est pas nécessairement sur la droite. L'intervalle de confiance sera plus grand dans le deuxième cas, pour tenir compte de la variabilité additionnelle causée par la fluctuation aléatoire autour de la droite.

\subsection{Prévision de la valeur moyenne $\esp[Y|x=x_0]$}

Soit $\esp[Y|x=x_0]$, la valeur moyenne de $Y$ lorsque $x=x_0$. Cette quantité est estimée par $\widehat{Y}_0 = \widehat{\beta}_0 + \widehat{\beta}_1 x_0$.
On utilise le fait que $Y_i \sim \N(\beta_0+\beta_1x_i,\sigma^2)$ pour déterminer la distribution de $\widehat{Y}_0$.
On exprime d'abord $\widehat{Y}_0$ sous la forme d'une combinaison linéaire des $Y_i$,
\begin{align*}
 \widehat{Y}_0 &= \widehat{\beta}_0 + \widehat{\beta}_1 x_0 \\
 &=\sum_{i=1}^n \left(\frac{1}{n} - a_i \overline{x}\right) Y_i + \left(\sum_{i=1}^n a_i Y_i\right) x_0 \\
 &= \sum_{i=1}^n \left(\frac{1}{n} - a_i \overline{x} + a_i x_0\right) Y_i\\
 &= \sum_{i=1}^n c_i Y_i,
 \intertext{où}
 c_i &= \frac{1}{n} - a_i \overline{x} + a_i x_0.
\end{align*}
On peut montrer que $\sum_{i=1}^n c_i = 1$, que $\sum_{i=1}^n c_i x_i = x_0$ et que $\sum_{i=1}^n c_i^2 = 1/n + (\overline{x}-x_0)^2 \big/ S_{xx}$.
On peut ensuite en déduire que
\begin{eqnarray}
 \esp\left(\widehat{Y}_0\right) &=& \sum_{i=1}^n c_i \esp\left(Y_i\right)= \sum_{i=1}^n c_i (\beta_0 + \beta_1 x_i) = \beta_0 + \beta_1 x_0, \\
 \label{eq:varyo}
 \var\left(\widehat{Y}_0\right) &=& \sum_{i=1}^n c_i^2 \var\left(Y_i\right) = \left\{ \frac{1}{n} + \frac{(\overline{x}-x_0)^2}{S_{xx}}\right\} \sigma^2.
 \end{eqnarray}
Puisque $\widehat{Y}_0$ est une combinaison linéaire des $Y_i$, qui suivent une loi normale, $\widehat{Y}_0$ suit aussi une loi normale. On a donc
\begin{equation}
\widehat{Y}_0 \sim \N\left\{\beta_0 + \beta_1 x_0, \left( 1/n + \frac{(\overline{x}-x_0)^2}{S_{xx}}\right) \sigma^2 \right\}.
\end{equation}
À partir de ce résultat, il est possible de calculer un intervalle de confiance autour de la valeur moyenne de $Y$ en $x=x_0$ en remplaçant $\sigma^2$ par son estimateur $\widehat\sigma^2$. En utilisant le même raisonnement qu'à la section \ref{sec:IntervallesDeConfiance}, on obtient l'intervalle de confiance de niveau $1-\alpha$ pour $\esp[Y|x=x_{0}]$,
\begin{eqnarray}
\left[\left(\widehat{\beta}_0 + \widehat{\beta}_1 x_0\right)\; \pm\; t_{n-2; 1-\alpha/2} \;\widehat\sigma \; \sqrt{\frac{1}{n}+\frac{(\overline x-x_0)^2}{S_{xx}}}\right].
\label{eq:int_conf_val_moy_prev}
\end{eqnarray}

 \subsection{Prévision d'une observation $Y_0^*$ lorsque $x=x_0$}

Il est important de comprendre que l'intervalle de confiance défini à l'équation \eqref{eq:int_conf_val_moy_prev} concerne la moyenne de $Y$ au point $x=x_0$. Pour construire ce qu'on appelle un intervalle de prévision, il faut considérer à la fois la variance liée à l'estimation des paramètres et la variance liée à la variabilité naturelle du phénomène, c'est-à-dire la variance du terme d'erreur $\varepsilon$. Or, l'équation \eqref{eq:varyo} ne tient pas compte de la variance du terme d'erreur $\sigma^2$.

Désignons la prévision de la moyenne de $Y$ en  $x_0$ par $\widehat{Y}_0 = \widehat{\beta}_0 + \widehat{\beta}_1 x_0$ et la prévision d'une observation individuelle par $\widehat{Y}_0^* = \widehat{\beta}_0 + \widehat{\beta}_1 x_0 + \varepsilon$. L'espérance de $\widehat{Y}_0$ et celle de $\widehat{Y}_0^*$ sont égales parce que l'espérance des termes d'erreur est nulle. Ce n'est pas le cas des variances.
On a que
\begin{eqnarray}
 \var\left(\widehat{Y}_0^*\right) &=& \var\left(\widehat{Y}_0\right) + \var \left(\varepsilon\right) \nonumber\\
 &=&\left\{ \frac{1}{n} + \frac{(\overline{x}-x_0)^2}{S_{xx}}\right\} \sigma^2 + \sigma^2 \nonumber\\
 &=& \left\{1+\frac{1}{n} + \frac{(\overline{x}-x_0)^2}{S_{xx}} \right\} \sigma^2.
 \label{eq:varyo*}
\end{eqnarray}
Le méthode de calcul de l'intervalle de prévision de niveau $1-\alpha$ pour ${Y}_0^*$ est similaire à celle de l'intervalle de confiance pour  $\esp[Y|x=x_0]$. On obtient donc
\begin{eqnarray}
\left[\left(\widehat{\beta}_0 + \widehat{\beta}_1 x_0\right)\; \pm\; t_{n-2; 1-\alpha/2} \; \widehat\sigma \; \sqrt{1+\frac{1}{n}+\frac{(\overline x-x_0)^2}{S_{xx}}} \right].
\label{eq:int_prev}
\end{eqnarray}


En analysant les expressions \eqref{eq:varyo} et \eqref{eq:varyo*}, nous constatons que les variances diminuent lorsque $n$ augmente ou que $S_{xx}$ augmente. Cela signifie que plus la taille de l'échantillon est grande et/ou plus il y a de la variabilité dans les valeurs de la variable indépendante, plus les estimations seront précises. Au contraire, plus le point $x_0$ pour lequel une prévision est requise est éloigné du centre  de la variable explicative $(\overline x)$, plus la prévision sera imprécise.
Il est important de prendre ces remarques en considération lors de la planification des expériences.


\begin{remarque}
L'intervalle de prévision pour $Y_0^*$ sera toujours plus large que l'intervalle de confiance pour $\esp[Y|x=x_0]$. Pour cette raison, il peut arriver dans la pratique qu'un client préfère présenter l'intervalle de confiance plutôt que l'intervalle de prévision. Il est important de bien comprendre et de bien expliquer la distinction entre les deux intervalles pour éviter une mauvaise interprétation des résultats.
\end{remarque}


\begin{exemple}
On désire calculer un intervalle de confiance au niveau 95 $\%$ pour le rendement moyen du produit ($\esp[Y|x=x_0]$) lorsque la température de fonctionnement est de 150 $^{\circ}\mathrm{C}$ et pour le rendement d'un produit spécifique (${Y}_0^*$) lorsque la température est de 150 $^{\circ}\mathrm{C}$.
Dans les deux cas, la prévision est $\widehat{Y}_0 = \widehat\beta_0+ \widehat\beta_1\times150 =-2,739+0,483\times 150=69,711 \%$.

La plupart des quantités nécessaires au calcul des deux intervalles ont été déjà calculées auparavant. On tire de la section \ref{exemple-RLS-interpretation}  que $\overline x=145, S_{xx}=8250$ et de l'exemple \ref{exp:int_beta} que $\widehat\sigma=0,950$ et que $t_{8; 0,975}=2,306$. Il nous reste seulement à calculer $(x_0-\overline x)^2=(150-145)^2=25$.
En utilisant l'équation \eqref{eq:int_conf_val_moy_prev}, on obtient l'intervalle de confiance du rendement moyen :
$$\left[69,711\pm 2,306\times 0,950\times\sqrt{1/10+25/8250}\right]=\left[69,711 \pm 0,703\right]=[69,008;70,414].$$
L'intervalle de prévision du rendement d'un produit quelconque, appelé également intervalle de prévision, se calcule à l'aide de l'équation~\eqref{eq:int_prev}: 
 $$\left[69,711\pm 2,306\times 0,950\times\sqrt{1+1/10+25/8250}\right]=[69,711\pm 2,301]=[67,410;72,012].$$
Comme mentionné, l'intervalle de prévision est plus large que l'intervalle de confiance pour $\esp[Y|x=150]$.
\end{exemple}


\section{Décomposition de la variance}
\label{sec:AnalyseDeLaVariance}




Une partie de la différence entre les $Y_i$ peut être expliquée par la différence entre les $x_i$, puisque selon le modèle de régression la variable réponse est expliquée en partie par la variable explicative. Il reste une partie de la variabilité des $Y_i$ qui n'est pas expliquée par les différences dans les valeurs de $x_i$. On peut écrire
\begin{align*}
Y_i - \overline{Y} &= \left(\widehat{Y}_i - \overline{Y}\right)  +  \left(Y_i - \widehat{Y}_i\right).
\end{align*}
On a aussi que
\begin{eqnarray*}
  \sum_{i=1}^n \left(Y_i-\overline{Y}\right)^2 & = & \sum_{i=1}^n \left(Y_i - \widehat{Y}_i + \widehat{Y}_i - \overline{Y}\right)^2\\
 &=&  \sum_{i=1}^n \left(Y_i - \widehat{Y}_i\right)^2 + \sum_{i=1}^n\left(\widehat{Y}_i - \overline{Y}\right)^2 + 2\sum_{i=1}^n \left(Y_i-\widehat{Y}_i\right)\left(\widehat{Y}_i-\overline{Y}\right)\\
 &=&  \sum_{i=1}^n \left(Y_i - \widehat{Y}_i\right)^2 + \sum_{i=1}^n\left(\widehat{Y}_i - \overline{Y}\right)^2.
\end{eqnarray*}
Notons
\begin{equation}\label{eq:decompvar}
SC_{tot}=SC_{reg}+SC_{res},
\end{equation}
o\`u
\begin{itemize}
    \item $SC_{tot} =\sum_{i=1}^n(Y_i-\overline{Y})^2$, la somme des carrés totale, quantifie la variabilit\'e des $Y_i$ autour de leur valeur moyenne;
    \item $SC_{reg} =\sum_{i=1}^n(\widehat{Y}_i-\overline{Y})^2$, la somme des carrés de la régression, quantifie la variabilit\'e des prévisions~$\widehat{Y}_i$. Les $\widehat{Y}_i$ ne varient qu'avec les $x_i$, cette somme de carrés est donc la partie de la variabilit\'e des $Y_i$ expliqu\'ee par le fait que toutes les observations n'ont pas la m\^eme valeur de $x_i$;
    \item $SC_{res} =\sum_{i=1}^n(Y_i-\widehat{Y}_i)^2$, la somme des carrés résiduelle, quantifie la variabilit\'e dans les $(Y_i-\widehat{Y}_i)$. Cette variabilit\'e est induite par le fait que la valeur de $Y_i$ n'est pas expliqu\'ee parfaitement par le mod\`ele de r\'egression.
\end{itemize}

\begin{exemple}
\label{exp:ss}
Toujours dans l'exemple du rendement d'un produit en fonction de la température de fonctionnement du procédé, on peut calculer 
\begin{align*}
SC_{tot} &= \sum_{i=1}^n(Y_i-\overline{Y})^2 = 1931,2, \\
SC_{res} &= \sum_{i=1}^n(Y_i-\widehat{Y}_i)^2 = (n-2)*\widehat\sigma^2 = 7,224, \\
SC_{reg} &= \sum_{i=1}^n(\widehat{Y}_i-\overline{Y})^2 = 1924,876.
\end{align*}
On vérifie alors que $SC_{tot} = SC_{reg} + SC_{res}$. On note que la température \`a elle seule ne peut pas expliquer en totalit\'e le fait que toutes les expériences ne possèdent pas le m\^eme rendement.
\end{exemple}

\subsection{Coefficient de détermination}
\label{sec:R2}


La décomposition de la variance permet de quantifier la qualité du modèle de régression linéaire à l'aide du coefficient de détermination,
\begin{equation}
R^2 =\frac{SC_{reg}} {SC_{tot}}.
\end{equation}
Notons que $0 \leq R^2 \leq 1$. Le coefficient de détermination s'interprète comme la proportion de la variance des $Y_i$ expliquée par le modèle. Lorsque $R^2=0$, le mod\`ele n'explique absolument pas la variabilité des $Y_i$. Quand $R^2=1$, tous les points sont alignés sur la droite de r\'egression, c'est-à-dire que l'ajustement du mod\`ele est parfait et que la valeur de $Y_i$ est une fonction exacte de $x_i$. Ainsi, une valeur élevée du $R^2$ indique que le modèle est bon.  Le seuil à partir duquel on considère que le $R^2$ est élevé est très variable et assez subjectif. Il dépend notamment des objectifs de la régression et du domaine d'application.

Le coefficient de détermination correspond au carré de la corrélation entre les valeurs observées et les valeurs prédites
 \[R^2 = \frac{SC_{reg}} {SC_{tot}} = \corr^2 (Y, \widehat{Y})=\left(\frac{\sum_{i=1}^n (Y_i - \overline{Y})(\widehat{Y}_i - \overline{Y})}{\sqrt{\sum_{i=1}^n (Y_i - \overline{Y})^2\sum_{i=1}^n(\widehat{Y}_i - \overline{Y})^2}}\right)^2.
\]


\begin{exemple}
Dans notre exemple,
$$R^2 = \frac{1924,876}{1931,1} = 0,9967,$$
ce qui indique que la température prédit presque parfaitement le rendement de l'expérience. Des valeurs aussi grandes sont plutôt rares en pratique.
\end{exemple}
